\begin{solution}{126}
    \begin{subsolution}
        如解题图2a所示，取球心为原点，过球心和球外点电荷所在处的直线为$X$轴（下同），导体球外点电荷$Q$位于$x = r \ ( r > R _ { 0 } )$，导体球外电荷及其像电荷位于$X$轴上（下同）。由接地导体球的球面电像公式知，点电荷$Q$的镜像电荷$Q _ { 1 } ^ { \prime }$的大小和位置分别为
        \begin{equation}
            Q _ {1} ^ {\prime} = - \frac {R _ {0}}{r} Q, \quad x ^ {\prime} = \frac {R _ {0} ^ {2}}{r}
            \label{eq:1.s126}
        \end{equation}
        由于不接地的导体球上电荷守恒以及导体球的外表面等势，在球心处还存在另一像电荷$Q _ { 2 } ^ { \prime }$，其大小和位置为
        \begin{equation}
            Q _ {2} ^ {\prime} = - Q _ {1} ^ {\prime} = \frac {R _ {0}}{r} Q, \quad x _ {2} ^ {\prime} = 0
            \label{eq:2.s126}
        \end{equation}
        导体球内、外的像电荷（包括像电偶极距）大小及其位置的坐标分别用打撇、不打撇加以区分（下同，特别声明的除外）。
    \end{subsolution}
    \begin{subsolution}
        设在$X$轴上有等量异号点电荷$Q _ { 1 } = Q \ ( Q > 0 )$和$Q _ { 2 } = - Q$，与金属球（导体球）球心$O$的距离均为$r \ ( r > R _ { 0 } )$，如解题图2b所示。由球面电像公式可知，$Q _ { 1 }$和$Q _ { 2 }$的镜像电荷$Q _ { 1 } ^ { \prime }$和$Q _ { 2 } ^ { \prime }$的大小和位置分别为
        \begin{equation}
            \left\{ \begin{array}{l} Q _ {1} ^ {\prime} = - \frac {R _ {0}}{r} Q _ {1} = - \frac {R _ {0}}{r} Q, \quad x _ {1} ^ {\prime} = - b \equiv - \frac {R _ {0} ^ {2}}{r}; \\ Q _ {2} ^ {\prime} = - \frac {R _ {0}}{r} Q _ {2} = \frac {R _ {0}}{r} Q, \quad x _ {2} ^ {\prime} = b \equiv \frac {R _ {0} ^ {2}}{r}. \end{array} \right.
            \label{eq:3.s126}
        \end{equation}
        像电荷$Q _ { 1 } ^ { \prime }$和$Q _ { 2 } ^ { \prime }$之和等于零，已保持导体球为电中性。当$r \to \infty$时，$b \to 0$；一对像电荷$Q _ { 1 } ^ { \prime }$和$Q _ { 2 } ^ { \prime }$可视为一位于球心$O$的点电偶极子，其偶极矩$p ^ { \prime }$的大小为
        \begin{equation}
            p ^ {\prime} = Q _ {2} ^ {\prime} \left(x _ {2} ^ {\prime} - x _ {1} ^ {\prime}\right) = 2 Q _ {1} ^ {\prime} b = 2 \frac {R _ {0} ^ {3}}{r ^ {2}} Q
            \label{eq:4.s126}
        \end{equation}
        匀强外电场$E _ { 0 }$可视为等量异号点电荷$\pm Q$在其连线中点处产生的电场，有
        \begin{equation}
            E _ {0} = 2 \frac {1}{4 \pi \varepsilon_ {0}} \frac {Q}{r ^ {2}}
            \label{eq:5.s126}
        \end{equation}
        由此可得
        \begin{equation}
            p ^ {\prime} = 4 \pi \varepsilon_ {0} E _ {0} R _ {0} ^ {3}
            \label{eq:6.s126}
        \end{equation}
    \end{subsolution}
    \begin{subsolution}
        由题设，球外点电荷$q _ { 1 } = - q \ ( q > 0 )$和$q _ { 2 } = q$分别位于$x _ { 1 } = r - \frac { \Delta l } { 2 }$和$x _ { 2 } = r + \frac { \Delta l } { 2 }$，如解题图2c所示。由\eqref{eq:1.s126}式知，$q _ { 1 }$和$q _ { 2 }$的像电荷$q _ { 1 } ^ { \prime }$和$q _ { 2 } ^ { \prime }$的大小和位置分别为
        \begin{align}
            q _ {1} ^ {\prime} &= \frac {q R _ {0}}{r - \frac {\Delta l}{2}} > 0, \quad x _ {1} ^ {\prime} = \frac {R _ {0} ^ {2}}{r - \frac {\Delta l}{2}} \label{eq:7.s126} \\
            q _ {2} ^ {\prime} &= - \frac {q R _ {0}}{r + \frac {\Delta l}{2}} < 0, \quad x _ {2} ^ {\prime} = \frac {R _ {0} ^ {2}}{r + \frac {\Delta l}{2}} \label{eq:8.s126}
        \end{align}
        由于导体球上电荷守恒和导体球的外表面等势，在球心处还存在另一像电荷$Q ^ { \prime }$，其大小和位置分别为
        \begin{equation}
            Q ^ {\prime} = - \left(q _ {1} ^ {\prime} + q _ {2} ^ {\prime}\right) = - \frac {R _ {0} q \Delta l}{r ^ {2} - \left(\frac {\Delta l}{2}\right) ^ {2}} = - \frac {R _ {0}}{r ^ {2}} q \Delta l < 0, \quad x _ {3} ^ {\prime} = 0
            \label{eq:9.s126}
        \end{equation}
        这里，已利用条件$\Delta l << r$（下同）。

        $q _ { 1 } ^ { \prime }$和$q _ { 2 } ^ { \prime }$组成的像电荷体系可视为位于
        \begin{equation}
            x _ {0} ^ {\prime} = b \equiv \frac {R _ {0} ^ {2}}{r}
            \label{eq:10.s126}
        \end{equation}
        点的像点电荷，电荷量为
        \begin{equation}
            - Q ^ {\prime} = q _ {1} ^ {\prime} + q _ {2} ^ {\prime} = \frac {R _ {0}}{r ^ {2}} q \Delta l = - q _ {3} ^ {\prime} > 0
            \label{eq:11.s126}
        \end{equation}
        和一个像点电偶极子，电偶极距大小为
        \begin{equation}
            p ^ {\prime} = \frac {R _ {0} ^ {3}}{r ^ {3}} q \Delta l
            \label{eq:12.s126}
        \end{equation}
        方向沿$X$轴正向，即与球外电偶极距的方向相同。

        此外，一对点电荷$Q ^ { \prime } < 0$，$- Q ^ { \prime } > 0$体系具有不为零的点电偶极距（以及高阶距），电偶极距的大小为
        \begin{equation}
            p ^ {\prime \prime} = (- Q ^ {\prime}) (x _ {0} ^ {\prime} - 0) = \frac {R _ {0} ^ {3}}{r ^ {3}} q \Delta l
            \label{eq:13.s126}
        \end{equation}
        方向沿$X$轴正向，即与球外电偶极距的方向相同。

        合起来，球外偶极子的镜像电偶极距的大小为
        \begin{equation}
            p _ {\text {总}} ^ {\prime} = p ^ {\prime} + p ^ {\prime \prime} = 2 \frac {R _ {0} ^ {3}}{r ^ {3}} q \Delta l
            \label{eq:14.s126}
        \end{equation}
        方向沿$X$轴正向，即与球外电偶极距的方向相同。

        [解法（二）

        设球外的电偶极子是点电荷$q _ { 1 } = - q \ ( q > 0 )$位于$x _ { 1 } = r - \frac { \Delta l } { 2 }$和$q _ { 2 } = q$位于$x _ { 2 } = r + \frac { \Delta l } { 2 }$，那么$q _ { 1 }$有两个像电荷，分别是
        \begin{align}
            q _ {1} ^ {\prime} &= \frac {q R _ {0}}{r - \frac {\Delta l}{2}} (> 0), \quad x _ {1} ^ {\prime} = \frac {R _ {0} ^ {2}}{r - \frac {\Delta l}{2}}; \label{eq:15.s126} \\
            q _ {1} ^ {\prime \prime} &= - \frac {q R _ {0}}{r - \frac {\Delta l}{2}} (< 0), \quad x _ {1} ^ {\prime \prime} = 0. \label{eq:16.s126}
        \end{align}
        同理，$q _ { 2 }$也有两个像电荷，分别是
        \begin{align}
            q _ {2} ^ {\prime} &= - \frac {q R _ {0}}{r + \frac {\Delta l}{2}} (< 0), \quad x _ {2} ^ {\prime} = \frac {R _ {0} ^ {2}}{r + \frac {\Delta l}{2}}; \label{eq:17.s126} \\
            q _ {2} ^ {\prime \prime} &= \frac {q R _ {0}}{r + \frac {\Delta l}{2}} (> 0), \quad x _ {2} ^ {\prime \prime} = 0. \label{eq:18.s126}
        \end{align}
        四个像电荷$q _ { 1 } ^ { \prime } , q _ { 1 } ^ { \prime \prime } , q _ { 2 } ^ { \prime } , q _ { 2 } ^ { \prime \prime }$构成的电荷系统的总电偶极矩为
        \begin{equation}
            p _ {\text {总}} ^ {\prime} = q _ {1} ^ {\prime} x _ {1} ^ {\prime} + q _ {1} ^ {\prime \prime} x _ {1} ^ {\prime \prime} + q _ {2} ^ {\prime} x _ {2} ^ {\prime} + q _ {2} ^ {\prime \prime} x _ {2} ^ {\prime \prime}
            \label{eq:19.s126}
        \end{equation}
        即
        \begin{equation}
            p _ {\text {总}} ^ {\prime} = \frac {q R _ {0} ^ {3}}{\left[ r - (\Delta l / 2) \right] ^ {2}} - \frac {q R _ {0} ^ {3}}{\left[ r + (\Delta l / 2) \right] ^ {2}} = 2 \frac {R _ {0} ^ {3}}{r ^ {3}} q \Delta l
            \label{eq:20.s126}
        \end{equation}
        其方向与球外电偶极距的方向相同。
        ]
    \end{subsolution}
    \begin{subsolution}
        此后，对于导体球内的像电荷（包括像电偶极距）大小及其位置的坐标取消打撇的标志。如图所示，当空间电场方向有两个相同导体球时，在恒外电场作用下，在两球心位置处产生一对点电偶极子
        \begin{equation}
            p _ {1} = 4 \pi \varepsilon_ {0} E _ {0} R _ {0} ^ {3}
            \label{eq:21.s126}
        \end{equation}
        这一对像偶极子又会产生次级电偶极子$p _ { 2 }$以及位于电偶极子$p _ { 2 }$处和球心的一对等量正反点电荷，接着产生$p _ { 3 } , p _ { 4 } , \cdots$及隔开的等量正反点电荷的镜像电荷（孤立点电荷），如此无穷地镜像映射下去（见解题图2e，其中孤立的镜像点电荷未画出）。由第（3）问分析知，各对点电偶极子的偶极矩和该偶极子所在位置到两球心的距离、以及孤立点电荷的大小依次为
        \begin{align}
            p _ {1} &= 4 \pi \varepsilon_ {0} E _ {0} R _ {0} ^ {3}, \quad b _ {1} = 0, \quad r _ {1} = r - b _ {1} = r \label{eq:22.s126} \\
            p _ {2} &= \frac {R _ {0} ^ {3}}{r _ {1} ^ {3}} p _ {1}, \quad b _ {2} = \frac {R _ {0} ^ {2}}{r _ {1}}, \quad r _ {2} = r - b _ {2}; \quad Q _ {1} = - q _ {1} ^ {(2)} < 0, \quad q _ {1} ^ {(2)} = \frac {R _ {0}}{r _ {1} ^ {2}} p _ {1} > 0 \label{eq:23.s126} \\
            p _ {3} &= \frac {R _ {0} ^ {3}}{r _ {2} ^ {3}} p _ {2}, \quad b _ {3} = \frac {R _ {0} ^ {2}}{r _ {2}}, \quad r _ {3} = r - b _ {3}; \label{eq:24.s126} \\
            Q _ {2} &= - \left(q _ {2} ^ {(2)} + q _ {2} ^ {(3)}\right), \quad q _ {2} ^ {(2)} = \frac {R _ {0}}{r _ {2}} Q _ {1} < 0, \quad q _ {2} ^ {(3)} = \frac {R _ {0}}{r _ {2}} \left(\frac {p _ {2}}{r _ {2}} + q _ {1} ^ {(2)}\right) > 0 \label{eq:25.s126} \\
            p _ {4} &= \frac {R _ {0} ^ {3}}{r _ {3} ^ {3}} p _ {3}, \quad b _ {4} = \frac {R _ {0} ^ {2}}{r _ {3}}, \quad r _ {4} = r - b _ {4}; \label{eq:26.s126} \\
            Q _ {3} &= - \left(q _ {3} ^ {(2)} + q _ {3} ^ {(3)} + q _ {3} ^ {(4)}\right), \quad q _ {3} ^ {(2)} = \frac {R _ {0}}{r _ {2}} Q _ {2}, \quad q _ {3} ^ {(3)} = \frac {R _ {0}}{r _ {3}} q _ {2} ^ {(2)}, \quad q _ {3} ^ {(4)} = \frac {R _ {0}}{r _ {3}} \left(\frac {p _ {3}}{r _ {3}} + q _ {2} ^ {(3)}\right) > 0 \label{eq:27.s126} \\
            p _ {n} &= \frac {R _ {0} ^ {3}}{r _ {n - 1} ^ {3}} p _ {n - 1}, \quad b _ {n} = \frac {R _ {0} ^ {2}}{r _ {n - 1}}, \quad r _ {n} = r - b _ {n}; \label{eq:28.s126} \\
            Q _ {n - 1} &= - \sum_ {i = 2} ^ {n} q _ {n - 1} ^ {(i)}, \quad q _ {n - 1} ^ {(2)} = \frac {R _ {0}}{r _ {2}} Q _ {n - 2} < 0, \quad q _ {n - 1} ^ {(3)} = \frac {R _ {0}}{r _ {3}} q _ {n - 2} ^ {(2)} < 0, \dots , \quad q _ {n - 1} ^ {(n - 1)} = \frac {R _ {0}}{r _ {n - 2}} q _ {n - 2} ^ {(n - 2)} < 0, \quad q _ {n - 1} ^ {(n)} = \frac {R _ {0}}{r _ {n - 1}} \left(\frac {p _ {n - 1}}{r _ {n - 1}} + q _ {n - 2} ^ {(n - 1)}\right) > 0 \label{eq:29.s126}
        \end{align}
        这里，$q _ { n } ^ { ( i ) }$是第$n$次镜像映射后左球内位置$b _ { i }$（即到球心的距离）上的镜像点电荷量，$Q _ { i }$是位于球心的补赏电荷（以保证导体球电中性以及球的外表面等势）。

        位于$\boldsymbol { b } _ { 0 }$、电偶极距为$\boldsymbol { p }$的电偶极子的电势为
        \begin{equation}
            \varphi (\boldsymbol {r}) = \frac {\boldsymbol {p} \cdot (\boldsymbol {r} - \boldsymbol {b} _ {0})}{4 \pi \varepsilon_ {0} | \boldsymbol {r} - \boldsymbol {b} _ {0} | ^ {3}}
            \label{eq:30.s126}
        \end{equation}
        取两球连线的中点为坐标原点，$x$轴水平向右，坐标原点为电势零点，则在球外空间任一点$P ( x , y )$处的电势为
        \begin{align}
            \varphi (x, y) &= - E _ {0} x + \sum_ {n = 1} ^ {\infty} \frac {(x + \frac {r}{2} - b _ {n}) p _ {n}}{4 \pi \varepsilon_ {0} [ (x + \frac {r}{2} - b _ {n}) ^ {2} + y ^ {2} ] ^ {\frac {3}{2}}} + \sum_ {n = 1} ^ {\infty} \frac {(x - \frac {r}{2} + b _ {n}) p _ {n}}{4 \pi \varepsilon_ {0} [ (x - \frac {r}{2} + b _ {n}) ^ {2} + y ^ {2} ] ^ {\frac {3}{2}}} \notag \\
            &\quad + \sum_ {n = 2} ^ {\infty} \sum_ {i = 2} ^ {n} \frac {q _ {n - 1} ^ {(i)}}{4 \pi \varepsilon_ {0} [ (x + \frac {r}{2} - b _ {i}) ^ {2} + y ^ {2} ] ^ {\frac {1}{2}}} + \sum_ {n = 2} ^ {\infty} \frac {Q _ {n - 1}}{4 \pi \varepsilon_ {0} [ (x + \frac {r}{2}) ^ {2} + y ^ {2} ] ^ {\frac {1}{2}}} \notag \\
            &\quad + \sum_ {n = 2} ^ {\infty} \sum_ {i = 2} ^ {n} \frac {q _ {n - 1} ^ {(i)}}{4 \pi \varepsilon_ {0} [ (x - \frac {r}{2} + b _ {i}) ^ {2} + y ^ {2} ] ^ {\frac {1}{2}}} + \sum_ {n = 2} ^ {\infty} \frac {Q _ {n - 1}}{4 \pi \varepsilon_ {0} [ (x - \frac {r}{2}) ^ {2} + y ^ {2} ] ^ {\frac {1}{2}}}
            \label{eq:31.s126}
        \end{align}
        当$y = 0$时，两球心连线上（球外）的电势分布为
        \begin{align}
            \varphi (x, y = 0) &= - E _ {0} x + \sum_ {n = 1} ^ {\infty} \frac {(x + \frac {r}{2} - b _ {n}) p _ {n}}{4 \pi \varepsilon_ {0} \left| x + \frac {r}{2} - b _ {n} \right| ^ {3}} + \sum_ {n = 1} ^ {\infty} \frac {(x - \frac {r}{2} + b _ {n}) p _ {n}}{4 \pi \varepsilon_ {0} \left| x - \frac {r}{2} + b _ {n} \right| ^ {3}} \notag \\
            &\quad + \sum_ {n = 2} ^ {\infty} \sum_ {i = 2} ^ {n} \frac {q _ {n - 1} ^ {(i)}}{4 \pi \varepsilon_ {0} \left| x + \frac {r}{2} - b _ {i} \right|} + \sum_ {n = 2} ^ {\infty} \frac {Q _ {n - 1}}{4 \pi \varepsilon_ {0} \left| x - \frac {r}{2} \right|} \notag \\
            &\quad + \sum_ {n = 2} ^ {\infty} \sum_ {i = 2} ^ {n} \frac {q _ {n - 1} ^ {(i)}}{4 \pi \varepsilon_ {0} \left| x + \frac {r}{2} - b _ {i} \right|} + \sum_ {n = 2} ^ {\infty} \frac {Q ^ {n - 1}}{4 \pi \varepsilon_ {0} \left| x - \frac {r}{2} \right|}
            \label{eq:32.s126}
        \end{align}
        两球心连线上（球外）的场强分布为
        \begin{align}
            E (x, y = 0) &= - \frac {\mathrm{d}}{\mathrm{d} x} \varphi (x, y = 0) = E _ {0} + \sum_ {n = 1} ^ {\infty} \frac {p _ {n}}{2 \pi \varepsilon_ {0} \left| x + \frac {r}{2} - b _ {n} \right| ^ {3}} + \sum_ {n = 1} ^ {\infty} \frac {p _ {n}}{2 \pi \varepsilon_ {0} \left| x - \frac {r}{2} + b _ {n} \right| ^ {3}} \notag \\
            &\quad + \sum_ {n = 2} ^ {\infty} \sum_ {i = 2} ^ {n} \frac {q _ {n - 1} ^ {(i)}}{4 \pi \varepsilon_ {0} \left| x + \frac {r}{2} - b _ {i} \right| ^ {2}} + \sum_ {n = 2} ^ {\infty} \frac {Q _ {n - 1}}{4 \pi \varepsilon_ {0} \left| x - \frac {r}{2} \right| ^ {2}} \notag \\
            &\quad + \sum_ {n = 2} ^ {\infty} \sum_ {i = 2} ^ {n} \frac {q _ {n - 1} ^ {(i)}}{4 \pi \varepsilon_ {0} \left| x + \frac {r}{2} - b _ {i} \right| ^ {2}} + \sum_ {n = 2} ^ {\infty} \frac {\mathcal {Q} _ {n - 1}}{4 \pi \varepsilon_ {0} \left| x - \frac {r}{2} \right| ^ {2}}
            \label{eq:33.s126}
        \end{align}
        式中，像电荷构成的电偶极子的方向与导体球外的外加匀强电场的方向相同。
    \end{subsolution}
    \begin{subsolution}
        设$r = 2 R _ { 0 }$并保持两球相互绝缘。由\eqref{eq:22.s126}、\eqref{eq:23.s126}、\eqref{eq:24.s126}式可知
        \begin{align}
            p _ {2} &= \left(\frac {1}{2}\right) ^ {3} p _ {1}, \quad b _ {2} = \frac {1}{2} R _ {0}, \quad r _ {2} = \frac {3}{2} R _ {0} \label{eq:34.s126} \\
            p _ {3} &= \left(\frac {2}{3}\right) ^ {3} \left(\frac {1}{2}\right) ^ {3} p _ {1} = \left(\frac {1}{3}\right) ^ {3} p _ {1}, \quad b _ {3} = \frac {2}{3} R _ {0}, \quad r _ {3} = \frac {4}{3} R _ {0} \label{eq:35.s126} \\
            &\vdots \notag \\
            p _ {n} &= \left(\frac {1}{n}\right) ^ {3} p _ {1}, \quad b _ {n} = \frac {n - 1}{n} R _ {0}, \quad r _ {n} = \frac {n + 1}{n} R _ {0}; \label{eq:36.s126}
        \end{align}
        由\eqref{eq:33.s126}式有
        \begin{align}
            E (x = 0, y = 0) &= E _ {0} + \sum_ {n = 1} ^ {\infty} \frac {p _ {n}}{\pi \varepsilon_ {0} \left| \frac {r}{2} - b _ {n} \right| ^ {3}} + \sum_ {n = 2} ^ {\infty} \sum_ {i = 2} ^ {n} \frac {q _ {n - 1} ^ {(i)}}{2 \pi \varepsilon_ {0} \left| \frac {r}{2} - b _ {i} \right| ^ {2}} + \sum_ {n = 2} ^ {\infty} \frac {Q _ {n}}{2 \pi \varepsilon_ {0} \left| \frac {r}{2} \right| ^ {2}} \notag \\
            &= E _ {0} + \sum_ {n = 1} ^ {\infty} \frac {p _ {n}}{\pi \varepsilon_ {0} \left| \frac {r}{2} - b _ {n} \right| ^ {3}} + \sum_ {n = 2} ^ {\infty} \sum_ {i = 2} ^ {n} \frac {q _ {n - 2} ^ {(i)}}{2 \pi \varepsilon_ {0} \left| \frac {r}{2} - b _ {i} \right| ^ {2}} + \sum_ {n = 2} ^ {\infty} \frac {- \sum_ {i = 2} ^ {n} q _ {n - 1} ^ {(i)}}{2 \pi \varepsilon_ {0} \left| \frac {r}{2} \right| ^ {2}} \notag \\
            &> E _ {0} + \sum_ {n = 1} ^ {\infty} \frac {p _ {n}}{\pi \varepsilon_ {0} \left| \frac {r}{2} - b _ {n} \right| ^ {3}} \notag \\
            &= E _ {0} + E _ {0} \sum_ {n = 1} ^ {\infty} \frac {4}{n ^ {3} \left| 1 - \frac {n - 1}{n} \right| ^ {3}} = E _ {0} + 4 E _ {0} \sum_ {n = 1} ^ {\infty} 1 \rightarrow \infty
            \label{eq:37.s126}
        \end{align}
        这里，用到了不等式
        \begin{equation}
            - Q _ {n - 1} \equiv \sum_ {i = 2} ^ {n} q _ {n - 1} ^ {(i)} > 0
            \label{eq:38.s126}
        \end{equation}
        这可通过直接计算证明如下：
        \begin{align}
            - Q _ {n - 1} &\equiv \sum_ {i = 2} ^ {n} q _ {n - 1} ^ {(i)} = \frac {R _ {0}}{r _ {2}} Q _ {n - 2} + \frac {R _ {0}}{r _ {3}} q _ {n - 2} ^ {(2)} + \dots + \frac {R _ {0}}{r _ {n - 2}} q _ {n - 2} ^ {(n - 1)} + \frac {R _ {0}}{r _ {n - 1}} q _ {n - 2} ^ {(n)} \notag \\
            &= \frac {R _ {0}}{r _ {2}} Q _ {n - 2} + \frac {R _ {0}}{r _ {3}} q _ {n - 2} ^ {(2)} + \dots + \frac {R _ {0}}{r _ {n - 2}} q _ {n - 2} ^ {(n - 1)} + \left[ \frac {R _ {0}}{r _ {n - 1}} \left(\frac {p _ {n - 1}}{r _ {n - 1}} + q _ {n - 2} ^ {(n - 1)}\right) \right] \notag \\
            &= \frac {R _ {0}}{r _ {2}} Q _ {n - 2} + \frac {R _ {0}}{r _ {3}} q _ {n - 2} ^ {(2)} + \dots + \frac {R _ {0}}{r _ {n - 2}} q _ {n - 2} ^ {(n - 2)} + \left[ \frac {R _ {0}}{r _ {n - 1} ^ {2}} p _ {n - 1} + \frac {R _ {0}}{r _ {n - 1}} \left(\frac {p _ {n - 2}}{r _ {n - 2}} + q _ {n - 3} ^ {(n - 2)}\right) \right] \notag \\
            &= \frac {R _ {0}}{r _ {2}} Q _ {n - 2} + \frac {R _ {0}}{r _ {3}} q _ {n - 2} ^ {(2)} + \dots + \frac {R _ {0}}{r _ {n - 2}} q _ {n - 2} ^ {(n - 2)} + \frac {R _ {0}}{r _ {n - 1}} q _ {n - 2} ^ {(n - 1)} + \frac {R _ {0}}{r _ {n - 1} ^ {2}} p _ {n - 1} \notag \\
            &> \frac {R _ {0}}{r _ {n - 1}} \left[ Q _ {n - 2} + \sum_ {i = 2} ^ {n - 1} q _ {n - 2} ^ {(i)} \right] + \frac {R _ {0}}{r _ {n - 1} ^ {2}} p _ {n - 1} = \frac {R _ {0}}{r _ {n - 1} ^ {2}} p _ {n - 1} > 0
            \label{eq:39.s126}
        \end{align}
    \end{subsolution}
\end{solution}








